% ==================================================================
%  徐森林、金亚东、薛春华《数学分析》（第三册）
%  复习题 13
%  源：原书印刷页 123-129（PDF 物理页 129-135）
%  底本：PaddleOCR-VL-1.6，见 chapters/_ch13_13_2_raw.md 末尾
%  约定：只用 exblock，不设 section（故不入目录）
%  全 27 题题面照录；订正说明见 tools/_make_ch13_ex.py 头部
% ==================================================================

\begin{exblock}{复习题 13}
\begin{enumerate}
  \item 设函数 $f(x)$ 在 $x=0$ 的邻域里有二阶连续导函数, 且 $f(0)=0$, $0<f'(0)<1$. 如果 $f_{n}(x)$ 为 $f(x)$ 的 $n$ 次复合, 即
  \[
    f_{n}=\underbrace{f\circ f\circ\cdots\circ f}_{n\text{ 次}}.
\]
  证明: 函数项级数 $\sum_{n=1}^{\infty}f_{n}(x)$ 在 $x=0$ 的邻域里一致收敛.
  \item 证明: 函数项级数
  \[
    \sum_{n=1}^{\infty}\frac{x^{n}}{1+x+x^{2}+\cdots+x^{2n-1}}\cos nx
\]
  在 $(0,1]$ 中一致收敛.
  \item 设 $\sum_{n=1}^{\infty}u_{n}(x)$ 在 $[a,b]$ 上收敛, 如果
  \[
    |S_{n}(x)|=\left|\sum_{k=1}^{n}u_{k}(x)\right|\leqslant M,\quad a\leqslant x\leqslant b,\ n=1,2,\cdots,
\]
  就称 $\sum_{n=1}^{\infty}u_{n}(x)$ 在 $[a,b]$ 上有界收敛.

  设 $\sum_{n=1}^{\infty}u_{n}(x)$ 在 $[a,b]$ 上有界收敛, 且对任意 $\delta>0$ 及 $c\in(a,b)$, $\sum_{n=1}^{\infty}u_{n}(x)$ 在 $[a,c-\delta]$ 与 $[c+\delta,b]$ 上一致收敛. 如果 $u_{n}(x)$ ($n=1,2,\cdots$) 在 $[a,b]$ 上 Riemann 可积, 则 $\sum_{n=1}^{\infty}u_{n}(x)$ 也在 $[a,b]$ 上 Riemann 可积, 且
  \[
    \int_{a}^{b}\left(\sum_{n=1}^{\infty}u_{n}(x)\right)\mathrm{d}x=\sum_{n=1}^{\infty}\int_{a}^{b}u_{n}(x)\mathrm{d}x.
\]
  \item 设函数项级数 $\sum_{n=1}^{\infty}u_{n}(x)$ 满足下列条件:

  (1) 广义积分 $\int_{a}^{+\infty}u_{n}(x)\mathrm{d}x$ 收敛, $n=1,2,\cdots$;

  (2) 函数项级数 $\sum_{n=1}^{\infty}u_{n}(x)$ 在区间 $[a,b]$ 上一致收敛, 这里 $b$ 是大于 $a$ 的任何实数;

  (3) 函数项级数 $\sum_{n=1}^{\infty}\int_{a}^{x}u_{n}(t)\mathrm{d}t$ 在 $[a,+\infty)$ 上一致收敛.

  则 $\int_{a}^{+\infty}\left(\sum_{n=1}^{\infty}u_{n}(x)\right)\mathrm{d}x$ 与 $\sum_{n=1}^{\infty}\int_{a}^{+\infty}u_{n}(x)\mathrm{d}x$ 都收敛, 且
  \[
    \int_{a}^{+\infty}\left(\sum_{n=1}^{\infty}u_{n}(x)\right)\mathrm{d}x=\sum_{n=1}^{\infty}\int_{a}^{+\infty}u_{n}(x)\mathrm{d}x.
\]
  \item 设 $r_{1},r_{2},\cdots$ 为 $[0,1]$ 中的有理数. 证明: 函数
  \[
    f(x)=\sum_{n=1}^{\infty}\frac{|x-r_{n}|}{3^{n}}\quad(0\leqslant x\leqslant1)
\]
  连续, 且在无理点可导, 但在有理点不可导.
  \item 设 $K(x,t)$ 在 $[a,b]\times[a,b]$ 上连续, $u_{0}(x)$ 在 $[a,b]$ 上连续, 且对任意 $x\in[a,b]$, 令
  \[
    u_{n}(x)=\int_{a}^{x}K(x,t)u_{n-1}(t)\mathrm{d}t,\quad n=1,2,\cdots,
\]
  证明: 函数列 $\{u_{n}(x)\}$ 在 $[a,b]$ 上一致收敛.
  \item 设 $f(x)$ 在 $(-\infty,+\infty)$ 内连续; 当 $x\neq0$ 时, 有 $|f(x)|<|x|$. 令
  \[
    \begin{aligned}
      f_{1}(x)&=f(x),\quad f_{2}(x)=f(f_{1}(x))=f\circ f(x),\\
      f_{n}(x)&=f(f_{n-1}(x))=\underbrace{f\circ f\circ\cdots\circ f}_{n\text{ 个}}(x),\ \cdots
    \end{aligned}
\]
  证明: $\{f_{n}(x)\}$ 在 $[-A,A]$ 上一致收敛.
  \item 设 $\alpha_{n}>0$, $\lim\limits_{n\to+\infty}\alpha_{n}=0$, $|u_{n}(x)|\leqslant\alpha_{n}$, $\forall x\in I$, 且 $u_{i}(x)u_{j}(x)=0$ ($i\neq j$). $\forall x\in I\subset\mathbb{R}$.

  证明: $\sum_{n=1}^{\infty}u_{n}(x)$ 在 $I$ 上一致收敛.
  \item 设 $\{u_{n}(x)\}$ 为 $[a,b]$ 上的可导函数列, 且在 $[a,b]$ 上有
  \[
    \left|\sum_{k=1}^{n}u_{k}'(x)\right|\leqslant C,
\]
  其中 $C$ 是不依赖 $x$ 与 $n$ 的正数. 证明: 若 $\sum_{n=1}^{\infty}u_{n}(x)$ 在 $[a,b]$ 上收敛, 则必为一致收敛.
  \item 设函数列 $f_{0}(x)$, $f_{1}(x)$, $\cdots$, $f_{n}(x)$, $\cdots$ 在 $I\subseteq\mathbb{R}$ 上有定义, 且满足:
  \[
    |f_{0}(x)|\leqslant M,\quad x\in I,
\]
  \[
    \sum_{n=0}^{m}|f_{n}(x)-f_{n+1}(x)|\leqslant M,\quad m=0,1,2,\cdots,\ x\in I,
\]
  其中 $M$ 为常数. 证明: 如果级数 $\sum_{n=0}^{\infty}b_{n}$ 收敛, 则函数项级数 $\sum_{n=0}^{\infty}b_{n}f_{n}(x)$ 必在 $I$ 上一致收敛.
  \item 证明: 函数项级数
  \[
    \sum_{n=1}^{\infty}\frac{1}{n}\left[\mathrm{e}^{x}-\left(1+\frac{x}{n}\right)^{n}\right]
\]
  在任意有限区间 $[a,b]$ 上一致收敛, 但在 $(0,+\infty)$ 上非一致收敛.
  \item 证明: $\sum_{n=1}^{\infty}(1-x)\frac{x^{n}}{1-x^{2n}}\sin nx$ 在 $\left(\frac{1}{2},1\right)$ 内一致收敛.
  \item 证明:

  (1) 函数列 $\left(1+\frac{x}{n}\right)^{n}$ ($n=1,2,\cdots$) 在 $[0,1]$ 上一致收敛于 $\mathrm{e}^{x}$;

  (2) 函数列 $f_{n}(x)=\dfrac{1}{\mathrm{e}^{\frac{x}{n}}+\left(1+\frac{x}{n}\right)^{n}}$ ($n=1,2,\cdots$) 在 $[0,1]$ 上一致收敛于 $\dfrac{1}{1+\mathrm{e}^{x}}$;

  (3) $\lim\limits_{n\to+\infty}\int_{0}^{1}\frac{\mathrm{d}x}{\mathrm{e}^{\frac{x}{n}}+\left(1+\frac{x}{n}\right)^{n}}=1+\ln\frac{2}{1+\mathrm{e}}$.
  \item 设 $\{r_{n}\}$ 为区间 $[0,1]$ 上全体有理数组成的数列, $u_{n}(x)$ 在 $[0,1]$ 上一致地满足 Lipschitz 条件, 即 $\exists L>0$, 使得 $\forall x_{1},x_{2}\in[0,1]$ 有
  \[
    \left|u_{n}(x_{1})-u_{n}(x_{2})\right|\leqslant L|x_{1}-x_{2}|,\ \forall n\in\mathbb{N}.
\]
  又设 $u_{n}(0)=0$, $u_{n}(x)$ 只在 $r_{n}$ 处不可导, 而在 $[0,1]$ 的其他点处可导 ($n=1,2,\cdots$). 证明:

  (1) $f(x)=\sum_{n=1}^{\infty}\frac{u_{n}(x)}{2^{n}}$ 在 $[0,1]$ 上连续;

  (2) $f(x)$ 在 $[0,1]$ 的有理点处不可导, 而在无理点处可导, 且导数
  \[
    f'(x)=\sum_{n=1}^{\infty}\frac{u_{n}'(x)}{2^{n}}.
\]
  \item 设 $g(x)$ 与 $f_{n}(x)\geqslant0$ ($n=1,2,\cdots$) 在 $[a,b]$ 上有界可积, 且 $\forall c\in(a,b)$, 有
  \[
    f_{n}(x)\rightrightarrows0,\ x\in[c,b],\quad n\to+\infty;
\]
  \[
    \lim_{n\to+\infty}\int_{a}^{b}f_{n}(x)\mathrm{d}x=1,\quad\lim_{x\to a^{+}}g(x)=A.
\]
  证明: $\lim\limits_{n\to+\infty}\int_{a}^{b}g(x)f_{n}(x)\mathrm{d}x=A$.
  \item 证明: 函数项级数 $\sum_{n=1}^{\infty}x^{2n}\ln x$ 在 $(0,1)$ 内不一致收敛, 但在 $[0,1]$ 上可逐项积分.
  \item 设 $\{f_{n}(x)\}$ 为 $(-\infty,+\infty)$ 上定义并且连续的函数列. 试构造一个在 $(-\infty,+\infty)$ 上有界、连续的函数 $f(x)$, 使得 $f(x)=0\Leftrightarrow f_{n}(x)=0$, $n=1,2,\cdots$.
  \item 设 $f_{n}(x)$ ($n=1,2,\cdots$) 在 $[0,1]$ 上连续, 并且 $f_{n}(x)\geqslant f_{n+1}(x)$, $\forall x\in[0,1]$, $n=1,2,\cdots$. 如果 $f_{n}(x)$ 在 $[0,1]$ 上收敛于 $f(x)$. 证明: $f(x)$ 在 $[0,1]$ 上达到最大值. 举例说明, 上述 $f(x)$ 在 $[0,1]$ 上未必达到最小值.
  \item 设 $f_{n}(x)$ ($n=1,2,\cdots$) 在 $[a,b]$ 上连续可导, 且 $\max\limits_{a\leqslant x\leqslant b}|f_{n}'(x)|\leqslant M$ ($n=1,2,\cdots$). 又 $\{f_{n}(x)\}$ 在 $[a,b]$ 上每一点处收敛. 证明: $\{f_{n}(x)\}$ 在 $[a,b]$ 上一致收敛. 如果 ``连续可导'' 改为 ``可导'', 上述结论是否仍正确?
  \item 设 $\{f_{n}(x)\}$ 是定义在 $[a,b]$ 上的连续函数列, 且 $\{f_{n}(x)\}$ 在 $[a,b]$ 上一致收敛于 $f(x)$. 又 $\{x_{n}\}\subset[a,b]$, 且 $x_{n}\to x_{0}$ ($n\to+\infty$), 证明: $\lim\limits_{n\to+\infty}f_{n}(x_{n})=f(x_{0})$.
  \item (Dini 定理) 设 $\{f_{n}(x)\}$ 为 $[a,b]$ 上的连续函数列, 且对任意 $x\in[a,b]$, 有 $f_{n}(x)\leqslant f_{n+1}(x)$. 证明: 如果 $\{f_{n}(x)\}$ 收敛于连续函数 $f(x)$, 则 $\{f_{n}(x)\}$ 在 $[a,b]$ 上必一致收敛于 $f(x)$.
  \item 设对每一个 $n$, $f_{n}(x)$ 为 $[a,b]$ 上的单调函数. 证明: 如果函数列 $\{f_{n}(x)\}$ 在 $[a,b]$ 上收敛于连续函数 $f(x)$, 则 $\{f_{n}(x)\}$ 在 $[a,b]$ 上必一致收敛于 $f(x)$.
  \item 设每个函数 $u_{n}(x)$ ($n=1,2,\cdots$) 都在点 $x=c$ 处连续, 但 $\sum_{n=1}^{\infty}u_{n}(c)$ 发散. 证明: $\forall\delta>0$, 级数 $\sum_{n=1}^{\infty}u_{n}(x)$ 在 $(c,c+\delta)$ 内不一致收敛.
  \item 设 $\{u_{n}(x)\}$ 是 $[a,b]$ 上正的单调减 ($u_{n+1}(x)\leqslant u_{n}(x)$) 且收敛于 $0$ 的函数列, 每个 $u_{n}(x)$ 都是 $[a,b]$ 上的单调增函数. 证明: 函数项级数 $\sum_{n=1}^{\infty}(-1)^{n-1}u_{n}(x)$ 在 $[a,b]$ 上一致收敛.
  \item 证明: 函数项级数 $\sum_{n=1}^{\infty}\frac{n^{2}}{\sqrt{n!}}(x^{n}+x^{-n})$ 在 $\frac{1}{2}\leqslant x\leqslant1$ 上一致收敛.
  \item 证明: Riemann $\zeta$ 函数 $\zeta(x)=\sum_{n=1}^{\infty}\frac{1}{n^{x}}$ 在 $(1,+\infty)$ 上为 $C^{\infty}$ 函数. 但它在 $(1,+\infty)$ 上不一致连续. 且 $\sum_{n=1}^{\infty}\frac{1}{n^{x}}$ 在 $(1,+\infty)$ 上非一致收敛.
  \item 设级数 $\sum_{n=1}^{\infty}a_{n}$ 收敛, 证明: Dirichlet 级数 $\sum_{n=1}^{\infty}\frac{a_{n}}{n^{x}}$ 当 $x\geqslant0$ 时一致收敛.
\end{enumerate}
\end{exblock}
